\section{Identifiability Discussion}
\label{sec:identifiability}

The completed evidence supports practical non-identifiability, not a claim of global
mathematical non-identifiability over all parameter space. Three independent lines are
relevant.

\paragraph{Development diagnostics.}
Earlier market-supported G2 studies found full algebraic rank in many cases while
failing target-blind recovery under clean and noisy observations. A bounded four-case
global study produced 40 clean near-equivalent solutions forming 39 complete-linkage
clusters; median repricing RMSE was $4.708\times10^{-8}$ while median range-scaled
parameter RMSE was 0.1485. These are historical diagnostics relative to R2, but they
motivate the retained finding.

\paragraph{Representation selection.}
R3 improved local conditioning and clean best-start recovery, yet lost under the frozen
realistic-noise rule. At the primary 0.5\% level, selected-R2 best-start displacement
was approximately 0.383 range-scaled RMSE while fitted-surface error was about 0.884\%
relative RMSE. The third expiry adds clean information but does not resolve ambiguity
at observation scale.

\paragraph{Known-truth benchmark.}
On the final synthetic test split, traditional calibration achieves essentially exact
repricing on every surface; $94.00\%$ have range-scaled parameter RMSE $\leq0.25$, so
$6.00\%$ remain above that coarse threshold. This lower-bounds rather than exhausts the
ambiguity problem.

Compensated directions are visible in aggregate metrics. Total initial variance
$v^s_0+v^f_0$ and total long-run variance $\theta_s+\theta_f$ are estimated much better
than individual factors, while kappa or half-life parameters remain difficult for all
methods. Canonical slow/fast ordering is a declared tie-breaker, not evidence that an
unordered factor swap is economically identified.

The correct scientific statement is therefore:
\begin{center}
\texttt{PRACTICAL\_NON\_IDENTIFIABILITY = RETAINED\_RESEARCH\_FINDING}.
\end{center}
